\ifajar
  \section{Latihan}
\else
  \section{Bab 7: Keamanan Aplikasi Web}
\fi

Seluruh file yang dipakai pada latihan ini tersedia di
\berkas{sources/chapter07/}, dihitung relatif terhadap akar repositori.
Latihan membutuhkan Docker dengan Compose versi 2 dan Python 3. Listing
penyiapan masuk ke direktori itu lebih dahulu, dan seluruh perintah sesudahnya
dijalankan dari sana.

Seluruh pemeriksaan pada latihan ini bersifat mengamati dan hanya diarahkan
ke aplikasi yang berjalan di komputer sendiri. Tidak ada satu pun perintah di
bab ini yang boleh diarahkan ke situs kampus maupun sistem milik orang lain.
Hasil tiap mahasiswa akan berbeda pada bagian waktu dan daftar paket, dan
yang dilaporkan adalah pola perbandingannya.

Penyiapannya dilakukan sekali, seperti pada Listing \ref{lst:c7-siapkan}.

\begin{lstlisting}[language=bash, caption={Menyiapkan basis data dan venv, memakai \berkas{sources/chapter07/siapkan-data.sh}}, label={lst:c7-siapkan}]
cd sources/chapter07
docker compose up -d
./siapkan-data.sh
python3 -m venv ../.venv
../.venv/bin/pip install "psycopg[binary]" ruff pip-audit
\end{lstlisting}

\ifajar
  \subsection*{Latihan 1: Membaca Header Keamanan}
\else
  \subsection*{Prosedur 1: Membaca Header Keamanan}
\fi

Latihan ini membuktikan bahwa header keamanan adalah pilihan yang harus
dibuat, bukan sesuatu yang datang sendiri.

\begin{enumerate}
  \item Nyalakan aplikasinya, lalu nyalakan juga server statis bawaan Python
        sebagai pembanding, seperti pada Listing \ref{lst:c7-lat1-jalan}.
  \item Baca header kedua server dengan \texttt{curl}, lalu isi \tabisi{tab:c7-lat1-hasil}. Tabel \ref{tab:c7-lat1-contoh}
        memperlihatkan contoh yang sudah terisi.
  \item Tambahkan satu header baru ke \berkas{sources/chapter07/kebijakan.json},
        misalnya \texttt{X-Frame-Options} bernilai \texttt{DENY}, lalu
        nyalakan ulang aplikasinya dan periksa kembali.
  \item Jawab tiga pertanyaan berikut. Pertama, header mana yang membatasi
        asal skrip yang boleh dimuat halaman? Kedua, mengapa header keamanan
        dipasang di satu fungsi dan bukan di tiap alamat? Ketiga, apa yang
        terjadi pada header baru tadi tanpa aplikasinya dinyalakan ulang?
\end{enumerate}

\begin{lstlisting}[language=bash, caption={Membaca header dua server, tersedia di \berkas{sources/chapter07/aplikasi.py}}, label={lst:c7-lat1-jalan}]
source ../.venv/bin/activate
SECRET_SESI=rahasia-latihan python aplikasi.py &
(cd statis && python -m http.server 8031 &)
curl -sI http://127.0.0.1:8030/ \
  | grep -iE "^(content-security|x-content|referrer|permissions)"
curl -sI http://127.0.0.1:8031/index.html \
  | grep -iE "^(content-security|x-content|referrer|permissions)"

# Keluarannya:
# Content-Security-Policy: default-src 'self'
# X-Content-Type-Options: nosniff
# Referrer-Policy: no-referrer
# Permissions-Policy: geolocation=(), camera=(), microphone=()
\end{lstlisting}

\begin{table}[!htb]
  \centering
  \caption{Contoh hasil Latihan 1 yang sudah terisi}
  \label{tab:c7-lat1-contoh}
  \begin{tabularx}{\textwidth}{Xll}
    \hline
    \textbf{Header} & \textbf{Server bawaan} & \textbf{Aplikasi Bab 7} \\
    \hline
    \texttt{Content-Security-}\newline\texttt{Policy} & tidak ada & \texttt{default-src 'self'} \\
    \hline
    \texttt{X-Content-Type-}\newline\texttt{Options} & tidak ada & \texttt{nosniff} \\
    \hline
  \end{tabularx}
\end{table}

\ifajar
\begin{table}[!htb]
  \centering
  \caption{Lembar isian hasil pengamatan Latihan 1}
  \label{tab:c7-lat1-hasil}
  \begin{tabularx}{\textwidth}{Xll}
    \hline
    \textbf{Header} & \textbf{Server bawaan} & \textbf{Aplikasi Bab 7} \\
    \hline
    \texttt{Content-Security-Policy} & \dots & \dots \\
    \hline
    \texttt{X-Content-Type-Options} & \dots & \dots \\
    \hline
    \texttt{Referrer-Policy} & \dots & \dots \\
    \hline
    \texttt{Permissions-Policy} & \dots & \dots \\
    \hline
    Header tambahan pilihan sendiri & \dots & \dots \\
    \hline
  \end{tabularx}
\end{table}
\fi

\ifajar
  \subsection*{Latihan 2: Cookie dan Aturan Lintas Origin}
\else
  \subsection*{Prosedur 2: Cookie dan Aturan Lintas Origin}
\fi

Latihan ini membuktikan bahwa izin lintas origin ditentukan daftar, bukan
kebetulan, dan bahwa atribut cookie terbaca langsung dari \textit{response}.

\begin{enumerate}
  \item Baca header \texttt{Set-Cookie} dari alamat profil, lalu catat tiap
        atributnya.
  \item Minta alamat yang sama tiga kali dengan header \texttt{Origin} yang
        berbeda, seperti pada Listing \ref{lst:c7-lat2-jalan}, lalu catat
        header izin yang kembali ke \tabisi{tab:c7-lat2-hasil}. Tabel
        \ref{tab:c7-lat2-contoh} memperlihatkan contoh yang sudah terisi.
  \item Tambahkan satu origin baru ke daftar di
        \berkas{sources/chapter07/kebijakan.json}, nyalakan ulang
        aplikasinya, lalu ulangi pemeriksaannya.
  \item Jawab tiga pertanyaan berikut. Pertama, siapa yang sebenarnya
        dilindungi aturan CORS? Kedua, mengapa aplikasi ini tidak memakai
        nilai bintang pada header izin? Ketiga, mengapa \texttt{Vary: Origin}
        tetap dikirim walaupun izinnya tidak diberikan?
\end{enumerate}

\begin{lstlisting}[language=bash, caption={Membaca cookie dan izin lintas origin, tersedia di \berkas{sources/chapter07/aplikasi.py}}, label={lst:c7-lat2-jalan}]
curl -sI http://127.0.0.1:8030/api/profil | grep -i "^set-cookie"
curl -sI -H "Origin: https://toko.pradita.ac.id" \
  http://127.0.0.1:8030/api/profil | grep -i "^access-control"
curl -sI -H "Origin: https://situs-lain.example" \
  http://127.0.0.1:8030/api/profil | grep -iE "^(access-control|vary)"

# Keluarannya:
# Set-Cookie: sesi=...; HttpOnly; SameSite=Lax; Path=/; Max-Age=3600
# Access-Control-Allow-Origin: https://toko.pradita.ac.id
# Vary: Origin
\end{lstlisting}

\begin{table}[!htb]
  \centering
  \caption{Contoh hasil Latihan 2 yang sudah terisi}
  \label{tab:c7-lat2-contoh}
  \begin{tabularx}{\textwidth}{lX}
    \hline
    \textbf{Origin yang dikirim} & \textbf{Header izin yang kembali} \\
    \hline
    \texttt{https://toko.pradita.ac.id} & berisi origin itu sendiri \\
    \hline
    \texttt{https://situs-lain.example} & tidak ada, hanya \texttt{Vary: Origin} \\
    \hline
  \end{tabularx}
\end{table}

\ifajar
\begin{table}[!htb]
  \centering
  \caption{Lembar isian hasil pengamatan Latihan 2}
  \label{tab:c7-lat2-hasil}
  \begin{tabularx}{\textwidth}{lX}
    \hline
    \textbf{Yang diperiksa} & \textbf{Hasilnya} \\
    \hline
    Atribut cookie sesi & \dots \\
    \hline
    Origin terdaftar pertama & \dots \\
    \hline
    Origin terdaftar kedua & \dots \\
    \hline
    Origin tidak terdaftar & \dots \\
    \hline
    Origin baru yang ditambahkan sendiri & \dots \\
    \hline
  \end{tabularx}
\end{table}
\fi

\ifajar
  \subsection*{Latihan 3: Teks Pengguna sebagai Data}
\else
  \subsection*{Prosedur 3: Teks Pengguna sebagai Data}
\fi

Latihan ini membuktikan bahwa teks bertanda baca tidak membuat basis data
maupun browser berpindah tugas.

\begin{enumerate}
  \item Jalankan pencarian biasa, lalu catat jumlah barisnya.
  \item Jalankan pencarian dengan kata kunci yang di dalamnya ada tanda kutip
        dan tanda titik koma, seperti pada Listing \ref{lst:c7-lat3-jalan},
        lalu catat jumlah baris dan status permintaannya.
  \item Buka halaman sapaan dengan nama yang berisi tanda kurung siku, lalu
        catat apa yang tertulis pada sumber halamannya.
\ifajar
  \item Isi \tabisi{tab:c7-lat3-hasil}. Tabel \ref{tab:c7-lat3-contoh}
        memperlihatkan contoh yang sudah terisi.
\else
  \item Bandingkan hasil pengukuran sendiri dengan Tabel
        \ref{tab:c7-lat3-contoh}.
\fi
  \item Jawab dua pertanyaan berikut. Pertama, mengapa pencarian kedua
        menghasilkan nol baris, dan apa yang akan terjadi bila teks itu
        disambung langsung ke SQL? Kedua, apa bedanya menolak teks tertentu
        dengan memperlakukan seluruh teks sebagai data?
\end{enumerate}

\begin{lstlisting}[language=bash, caption={Mengirim teks bertanda baca sebagai kata kunci, tersedia di \berkas{sources/chapter07/aplikasi.py}}, label={lst:c7-lat3-jalan}]
curl -s "http://127.0.0.1:8030/api/cari?judul=Pengantar%20Jilid%207" \
  | head -c 120
curl -s --get --data-urlencode "judul=Pengantar'; SELECT 1 --" \
  http://127.0.0.1:8030/api/cari | head -c 120
curl -s --get --data-urlencode 'nama=<b>Mahasiswa</b>' \
  http://127.0.0.1:8030/sapa | grep -i halo

# Keluarannya:
# {"kata_kunci": "Pengantar Jilid 7", "hasil": [{"id": ...
# {"kata_kunci": "Pengantar'; SELECT 1 --", "hasil": []}
# <body><h1>Halo, &lt;b&gt;Mahasiswa&lt;/b&gt;</h1>
\end{lstlisting}

\begin{table}[!htb]
  \centering
  \caption{Contoh hasil Latihan 3 yang sudah terisi}
  \label{tab:c7-lat3-contoh}
  \begin{tabularx}{\textwidth}{Xrl}
    \hline
    \textbf{Yang dikirim} & \textbf{Baris} & \textbf{Status} \\
    \hline
    \texttt{Pengantar Jilid 7} & 20 & 200 \\
    \hline
    \texttt{Pengantar'; SELECT 1 -{}-} & 0 & 200 \\
    \hline
  \end{tabularx}
\end{table}

\ifajar
\begin{table}[!htb]
  \centering
  \caption{Lembar isian hasil pengamatan Latihan 3}
  \label{tab:c7-lat3-hasil}
  \begin{tabularx}{\textwidth}{Xrl}
    \hline
    \textbf{Yang dikirim} & \textbf{Baris} & \textbf{Status} \\
    \hline
    Kata kunci biasa & \dots & \dots \\
    \hline
    Kata kunci bertanda baca & \dots & \dots \\
    \hline
    Nama berisi tanda kurung siku & \dots & \dots \\
    \hline
  \end{tabularx}
\end{table}
\fi

\ifajar
  \subsection*{Latihan 4: Menegakkan Gate Keamanan}
\else
  \subsection*{Prosedur 4: Menegakkan Gate Keamanan}
\fi

Latihan ini memasang seluruh syarat di atas sebagai satu \textit{gate}, lalu
membandingkan \textit{verdict}-nya pada dua sasaran.

\begin{enumerate}
  \item Baca \berkas{sources/chapter07/kebijakan.json}, lalu tulis dugaan
        berapa syarat yang akan terpenuhi aplikasi dan berapa pada server
        pembanding.
  \item Jalankan \textit{gate}-nya terhadap keduanya seperti pada Listing
        \ref{lst:c7-lat4-jalan}, lalu catat jumlah syarat, jumlah yang lolos,
        dan status keluarnya ke \tabisi{tab:c7-lat4-hasil}. Tabel
        \ref{tab:c7-lat4-contoh} memperlihatkan contoh yang sudah terisi.
  \item Jalankan pemeriksaan dependensi dan analisis statis, lalu catat
        jumlah temuannya.
  \item Hentikan aplikasinya, lalu nyalakan ulang dengan \textit{copy} kebijakan
        yang kehilangan satu header, sementara \textit{gate}-nya tetap
        memakai \berkas{sources/chapter07/kebijakan.json} yang asli, seperti
        pada Listing \ref{lst:c7-lat4-hapus}. Sesudah hasilnya dicatat,
        nyalakan kembali aplikasinya tanpa argumen.
  \item Jawab tiga pertanyaan berikut. Pertama, mengapa \textit{gate} menilai
        \textit{response} yang dikirim dan bukan membaca kodenya? Kedua,
        mengapa satu syarat justru lolos pada server pembanding? Ketiga,
        mengapa hasil pemeriksaan dependensi punya umur simpan?
\end{enumerate}

\begin{lstlisting}[language=bash, caption={Menjalankan gate pada dua sasaran, tersedia di \berkas{sources/chapter07/periksa-keamanan.py}}, label={lst:c7-lat4-jalan}]
source ../.venv/bin/activate
python periksa-keamanan.py
echo $?
python periksa-keamanan.py kebijakan-pembanding.json
echo $?
./audit-dependensi.sh
ruff check .

# Keluarannya:
# Lolos: seluruh syarat keamanan terpenuhi.
# 0
# Gagal: 8 syarat keamanan belum terpenuhi.
# 1
# No known vulnerabilities found
# All checks passed!
\end{lstlisting}

\begin{lstlisting}[language=bash, caption={Aplikasi yang kehilangan satu header, memakai \berkas{sources/chapter07/aplikasi.py}}, label={lst:c7-lat4-hapus}]
grep -v '"Referrer-Policy"' kebijakan.json > kebijakan-uji.json
SECRET_SESI=rahasia-latihan python aplikasi.py kebijakan-uji.json &
python periksa-keamanan.py
echo $?

# Keluarannya (dipotong):
#   / Referrer-Policy                            tidak ada    LANGGAR
#   /api/cari?judul=Pengantar Referrer-Policy    tidak ada    LANGGAR
#   /api/profil Referrer-Policy                  tidak ada    LANGGAR
#
# Gagal: 3 syarat keamanan belum terpenuhi.
# 1
\end{lstlisting}

\begin{table}[!htb]
  \centering
  \caption{Contoh hasil Latihan 4 yang sudah terisi}
  \label{tab:c7-lat4-contoh}
  \begin{tabularx}{\textwidth}{Xrrr}
    \hline
    \textbf{Sasaran} & \textbf{Syarat} & \textbf{Lolos} & \textbf{Status keluar} \\
    \hline
    Aplikasi Bab 7 & 18 & 18 & 0 \\
    \hline
    Server statis bawaan Python & 9 & 1 & 1 \\
    \hline
    Aplikasi dengan satu header dihapus & 18 & 15 & 1 \\
    \hline
  \end{tabularx}
\end{table}

\ifajar
\begin{table}[!htb]
  \centering
  \caption{Lembar isian hasil pengamatan Latihan 4}
  \label{tab:c7-lat4-hasil}
  \begin{tabularx}{\textwidth}{Xrrr}
    \hline
    \textbf{Sasaran} & \textbf{Syarat} & \textbf{Lolos} & \textbf{Status keluar} \\
    \hline
    Aplikasi Bab 7 & \dots & \dots & \dots \\
    \hline
    Server statis bawaan Python & \dots & \dots & \dots \\
    \hline
    Aplikasi dengan satu header dihapus & \dots & \dots & \dots \\
    \hline
    Temuan dependensi dan analisis statis & \dots & \dots & \dots \\
    \hline
  \end{tabularx}
\end{table}
\fi

Setelah seluruh latihan selesai, container dihentikan dengan
\texttt{docker compose down -v}. Opsi \texttt{-v} ikut menghapus volume
datanya, sehingga penyiapan berikutnya dimulai dari keadaan bersih.
